\newpage
\section{Conclusion}
\label{sec:conclusion}

\subsection{Summary}
\label{subsec:conclusion-summary}

This report investigated joint task offloading and PRB allocation in a three-tier LOCAL--EDGE--CLOUD architecture for 6G-oriented edge-cloud computing. The problem was formulated as a Dec-POMDP and addressed through a CTDE-based MAPPO scheduler with a dual-head action design, a PPO comparison baseline, and a customized PureEdgeSim environment extended with PRB-aware networking, device-level observations, and Java--Python co-simulation.

The experimental results show a clear advantage of on-policy reinforcement learning over static heuristics. In the offline all-algorithm comparison, MAPPO achieves a task success rate of \(99.96\%\), PPO reaches \(99.94\%\), and the strongest heuristic baseline, ROUND\_ROBIN, remains at \(90.39\%\). Under the standard three-seed evaluation, MAPPO attains a mean success rate of \(99.9237\%\), which indicates that the proposed multi-agent scheduler is not merely feasible in this setting, but already stable and technically credible. The ablation study further isolates the source of this performance: removing agent embeddings causes only a small reduction to \(99.8148\%\), whereas fixing PRB allocation reduces mean success to \(92.3822\%\), showing that adaptive wireless-resource control is the dominant contributor to the final policy quality. The device-count study reaches the same systems-level conclusion from another angle: MAPPO maintains high success rates from 80 to 240 devices, while the deterioration at 320 devices is better explained by infrastructure saturation than by an ordinary algorithmic breakdown.

Taken together, the strongest conclusion of this report is that joint offloading and adaptive wireless-resource control is an effective design principle for this problem class. Within that joint-control framework, MAPPO provides the central methodological contribution of the report: it delivers reliable multi-agent coordination, remains competitive with a very strong PPO baseline, and offers a more convincing solution for deadline-sensitive tail control in the studied LOCAL--EDGE--CLOUD system.

\subsection{Limitations}
\label{subsec:conclusion-limitations}

Several limitations define the current boundary of this study. First, the evidence remains simulation-based, so a non-trivial gap still separates the modeled PRB-aware wireless environment from a real 5G or 6G deployment. Second, the present formulation studies full-task offloading only; it does not yet address partial offloading, where the scheduler would need to coordinate task partitioning together with destination selection and communication-resource allocation. Third, the evaluated architecture is limited to the LOCAL--EDGE--CLOUD setting and therefore does not yet cover a richer continuum in which neighboring devices or an explicit mist layer can also participate in execution.

Methodologically, the current MAPPO design should also be viewed as a strong but not exhaustive point in the design space. The critic input, reward coefficients, agent representation, and training-stability strategy were chosen carefully for this report, but a broader search over these factors would strengthen the claims further. Robustness validation also remains limited relative to the full ambition of the problem: although the preserved seeds and scaling study show stable behavior within the tested regime, broader verification across more seeds, traffic patterns, and topology perturbations is still needed. Finally, the finalized PPO artifacts do not preserve the same per-decision trajectory exports as MAPPO, so the aggregate PPO metrics are trustworthy but the behavioral comparison is not yet fully symmetric at the action level.

\subsection{Future Work}
\label{subsec:conclusion-future}

The current results suggest several direct directions for future work.
\begin{itemize}
    \item Extend the current full-offloading formulation to partial offloading, so that the scheduler jointly decides how a task should be partitioned, where each part should execute, and how the associated communication resources should be allocated.
    \item Expand the architecture beyond LOCAL--EDGE--CLOUD by introducing an explicit mist layer or neighboring-device execution, allowing tasks to be offloaded to nearby devices and thereby moving toward a fuller computing-continuum formulation.
    \item Further optimize MAPPO itself through a more systematic study of reward design, critic-state construction, agent representation, training-stability techniques, and checkpoint-selection strategy.
    \item Broaden the evaluation regime to denser deployments, alternative infrastructure capacities, more diverse traffic mixes, and stronger topology perturbations in order to test robustness beyond the workload range examined in this report.
    \item Reduce the simulation-to-reality gap by coupling the current framework with more realistic traffic traces or a small-scale hardware testbed, and by preserving equally detailed behavioral artifacts for all learned baselines.
\end{itemize}

Overall, the results support on-policy reinforcement learning as a practical approach to coordinated offloading and radio-resource control in hierarchical edge-cloud systems. More specifically, they show that adaptive PRB control is the key design factor behind the scheduler's quality, while MAPPO turns that joint-control principle into a credible multi-agent solution for deadline-sensitive edge-cloud orchestration. PPO remains a very strong benchmark, but the present evidence still places MAPPO at the center of the report's contribution because it combines near-best aggregate performance with a more persuasive multi-agent coordination story in the studied setting.
